\documentclass[12pt,a4paper]{article}

% Pacotes básicos
\usepackage[utf8]{inputenc}
\usepackage[T1]{fontenc}
\usepackage[brazil]{babel}
\usepackage{geometry}
\usepackage{graphicx}
\usepackage{amsmath, amssymb}
\usepackage{setspace}
\usepackage{hyperref}

% Margens
\geometry{
    left=3cm,
    right=2cm,
    top=3cm,
    bottom=2cm
}

% Espaçamento entre linhas
\onehalfspacing

% Informações do documento
\title{Título do Relatório}
\author{Pedro Henrique dos Reis Arcoverde- 254719}
\date{\today}

\begin{document}

% Capa
\begin{titlepage}
	\centering
	{\large UNICAMP}\\[1.5cm]

	{\Large MC 906}\\[2cm]

	{\Huge \textbf{PROJETO 2 - PORTIFÓLIO}}\\[2cm]

	{\large Pedro Henrique dos Reis Arcoverde-254719}\\[0.5cm]
	{\large Professor: Anderson}\\[2cm]

	{\large Cidade}\\
	{\large \today}
\end{titlepage}

% Resumo
\begin{abstract}
	Este é o resumo do relatório. Aqui você descreve brevemente o objetivo, metodologia e resultados.
\end{abstract}

% Sumário
\tableofcontents
\newpage

\section{Introdução}
Apresente o tema do relatório, contexto e objetivos.

\section{Fundamentação Teórica}
Explique os conceitos teóricos relevantes.

\subsection{Subseção Exemplo}
Detalhamento de algum tópico específico.

\section{Metodologia}
Descreva como o trabalho foi realizado.

\section{Experimentos}
Para esta seção vamos começar a olhar para uma implementação real dos nosso modelos de inteligência artificial:
Bellman e Q-Learning destrinchando as suas composições e alterando os valores das variáveis que regem o comportamento
dois dois de forma a tentar extrair o máximo de informção sobre o funcionamento, estratégia e rendimentos no fluxo de
decisão.

\subsection{Implementação das Equações de Bellman}
Utilizou-se o algoritmo de \textit{Value Iteration} para determinar a política ótima através da
atualização iterativa dos valores de estado $V(s)$:
\[ V(s) \leftarrow \max_{a} \sum_{s'} T(s,a,s') [R(s,a,s') + \gamma V(s')] \]
O experimento foi configurado com um fator de desconto $\gamma = 0.99$ e um critério de convergência $\theta = 1e-5$.
\subsection{Análise de Resultados: Value Iteration}
O algoritmo convergiu após \textbf{1.796 iterações}. Os resultados detalhados encontram-se na Tabela \ref{tab:v_table}.

\begin{table}[h]
	\centering
	\caption{Valores de Estado ($V(s)$) e Política Aprendida}
	\label{tab:v_table}
	\begin{tabular}{l|c|c}
		\hline
		\textbf{Estado}   & \textbf{Valor $V(s)$} & \textbf{Ação Ótima} \\ \hline
		0 (DOWN, CASH)    & 1,38E+09              & BUY                 \\
		1 (DOWN, HOLDING) & 1,41E+09              & HOLD                \\
		2 (FLAT, CASH)    & 1,34E+09              & SELL*               \\
		3 (FLAT, HOLDING) & 1,36E+09              & HOLD                \\
		4 (UP, CASH)      & 1,33E+09              & BUY                 \\
		5 (UP, HOLDING)   & 1,32E+09              & SELL                \\ \hline
	\end{tabular}
\end{table}

\subsection{Discussão Técnica}
Observou-se uma magnitude extremamente elevada nos valores da \textit{V-Table} (na ordem de $10^9$). Este fenômeno é
decorrente da combinação de um horizonte de planejamento longo com um fator de desconto $\gamma$ muito próximo
da unidade em um ambiente com tendência de crescimento (\textit{drift}) positiva.

\textbf{Análise Crítica:} A política ótima no Estado 2 indicou a ação de venda mesmo sem a posse do ativo
(\textit{CASH}). Tecnicamente, a recompensa esperada a longo prazo tornou-se tão vasta que a penalidade estática por
ação inválida ($r = -1.0$) tornou-se insignificante no cálculo do valor esperado.
Para mitigar este comportamento e garantir a estabilidade da política, recomenda-se a normalização das recompensas
ou a aplicação de penalidades proporcionais ao patrimônio líquido.

\subsection{Análise Experimental: Q-Learning}

Para avaliar o aprendizado \textit{model-free} e mitigar as a nomalias observadas no planejamento de Bellman,
foi conduzida uma busca em grade (\textit{Grid Search}) com o algoritmo Q-Learning. A análise a seguir
foca na configuração de melhor interpretabilidade e estabilidade, obtida com taxa de aprendizado
$\alpha = 0.1$ e fator de desconto $\gamma = 0.9$.

\subsubsection{Correção do Comportamento Anômalo (Ações Inválidas)}
No algoritmo de \textit{Value Iteration}, a combinação de horizonte infinito e tendência positiva (\textit{drift})
ofuscou a penalidade estática do ambiente ($r = -1.0$), levando o agente a tentar vender ativos que não possuía.
O Q-Learning corrigiu este comportamento com sucesso.

A inspeção da Tabela Q final revela que os valores esperados para ações inválidas convergiram quase perfeitamente
para a magnitude da penalidade configurada. Por exemplo, no Estado 0 (\textit{DOWN, CASH}), o valor da ação \textit{SELL}
atingiu $Q(s_0, a_2) \approx -0.966$. O mesmo ocorreu nos Estados 2 e 4. Isso comprova empíricamente que o agente
aprendeu as restrições físicas do MDP interagindo com o ambiente, evitando ações punitivas.

\subsection{Sensibilidade de Parâmetros no Q-Learning}

Para compreender a dinâmica de aprendizagem do agente em ambiente estocástico, foi executada uma pesquisa em grelha
(\textit{Grid Search}) avaliando 9 combinações de hiperparâmetros, cruzando a Taxa de Aprendizagem
($\alpha \in \{0.1, 0.5, 0.9\}$) e o Fator de Desconto ($\gamma \in \{0.5, 0.9, 0.99\}$). O treino consistiu em 2000 episódios
por configuração.

A Tabela \ref{tab:grid_search} sumariza
o estado final do agente (Episódio 2000) para cada combinação, destacando a Recompensa do episódio, o Saldo Final do portefólio
(capital inicial de \$1000.00) e a métrica de convergência $\Delta Q$ (diferença máxima absoluta na Matriz Q face à medição anterior).
\begin{table}[h]
	\centering
	\caption{Resultados Finais do \textit{Grid Search} de Q-Learning (Episódio 2000)}
	\label{tab:grid_search}
	\begin{tabular}{c c | c c c}
		\hline
		\textbf{$\alpha$ (Alpha)} & \textbf{$\gamma$ (Gamma)} & \textbf{Recompensa} & \textbf{Saldo Final (\$)} & \textbf{$\Delta Q$ (Convergência)} \\ \hline
		0.1                       & 0.50                      & -5.65               & 999.35                    & 0.5740                             \\
		0.1                       & 0.90                      & -3.51               & 998.49                    & 0.2781                             \\
		\textbf{0.1}              & \textbf{0.99}             & \textbf{2.32}       & \textbf{1002.32}          & \textbf{0.7140}                    \\ \hline
		0.5                       & 0.50                      & -2.11               & 999.89                    & 0.7805                             \\
		0.5                       & 0.90                      & -1.56               & 1003.44                   & 1.0843                             \\
		0.5                       & 0.99                      & 0.45                & 1000.45                   & 1.1782                             \\ \hline
		0.9                       & 0.50                      & -3.00               & 1000.00                   & 2.2827                             \\
		0.9                       & 0.90                      & -3.00               & 1000.00                   & 2.5645                             \\
		0.9                       & 0.99                      & -9.38               & 999.62                    & 3.4039                             \\ \hline
	\end{tabular}
\end{table}

\subsubsection{Interpretação dos Parâmetros}

\paragraph{O Impacto da Taxa de Aprendizagem ($\alpha$):} A métrica $\Delta Q$ revelou-se fundamental para aferir a estabilidade cognitiva do modelo.
Observou-se que taxas elevadas ($\alpha = 0.9$) são tóxicas para a alocação de portefólios: o $\Delta Q$ final manteve-se
criticamente elevado (entre 2.28 e 3.40), evidenciando um fenómeno de "amnésia".
O agente reagia de forma histérica ao ruído estocástico diário do mercado, reescrevendo a sua estratégia a cada flutuação.
Em contrapartida, com $\alpha = 0.1$, o $\Delta Q$ decaiu progressivamente para valores abaixo de 0.8. Esta configuração conferiu
inércia ao agente, permitindo-lhe filtrar o ruído e extrair o verdadeiro sinal das tendências do mercado.

\paragraph{O Impacto do Fator de Desconto ($\gamma$):} Fixando a taxa de aprendizagem ótima ($\alpha = 0.1$), o fator de desconto
ditou o perfil de risco do investidor. Com $\gamma = 0.5$, o agente adotou uma postura míope (\textit{high-frequency trading}),
empatando com o mercado por realizar lucros demasiado cedo. Com $\gamma = 0.9$, obteve-se o modelo mais estável matematicamente
($\Delta Q = 0.2781$), traduzindo-se num perfil hiperconservador focado apenas em não perder dinheiro. Já a configuração $\gamma = 0.99$
forçou o agente a comportar-se como um investidor institucional de longo prazo, suportando quedas temporárias para capturar os grandes
ciclos de alta.

\subsubsection{Seleção do Modelo Ótimo: $\alpha=0.1$ e $\gamma=0.99$}

A combinação de $\alpha=0.1$ e $\gamma=0.99$ emergiu inequivocamente como o modelo vencedor. A análise da sua evolução temporal
demonstrou três fases distintas:
\begin{enumerate}
	\item \textbf{Exploração Punitiva (Episódios 0 - 800):} O agente obteve recompensas severamente negativas e saldos abaixo dos \$1000,
	      resultado direto de testar aleatoriamente os limites do MDP (acionando multas por ações inválidas, como vender a descoberto).
	\item \textbf{Estabilização Cognitiva (Episódio 1200):} O $\Delta Q$ atingiu o seu mínimo histórico ($0.4052$), sinalizando matematicamente que o agente compreendeu as regras de transição estocástica do ambiente.
	\item \textbf{Explotação Lucrativa (Episódios 1600 - 2000):} Com a consolidação da Matriz Q e a redução do fator $\epsilon$, as recompensas
	      assumiram valores positivos consistentes (+4.88, +2.32), ultrapassando a barreira do capital inicial e encerrando o treino
	      em \$1002.32.
\end{enumerate}

\subsubsection{Conclusão}

Perante a aparente modéstia do lucro final (+\$2.32 após 2000 episódios), é imperativo contextualizar a natureza da simulação
quantitativa. O ambiente foi deliberadamente modelado através de um \textit{Random Walk} Geométrico com um forte peso estocástico,
complementado por restrições operacionais rígidas e penalizações financeiras severas.

Neste contexto, o sucesso da Inteligência Artificial não reside em memorizar um histórico de preços para gerar um retorno irrealista
de 100\% (o que indicaria \textit{overfitting}), mas sim em aprender a operar estatisticamente. A vitória estrutural deste modelo
($\alpha=0.1, \gamma=0.99$) é a sua capacidade de sobreviver. O agente conseguiu estancar a hemorragia de perdas resultantes da
ignorância inicial, mitigou o erro de diferença temporal ($\Delta Q$), adaptou-se a uma dinâmica ruidosa e concluiu milhares
de ciclos com o património líquido protegido e com lucro estatisticamente robusto, cumprindo em pleno os requisitos da formulação
do Processo de Decisão de Markov.
\section{Conclusão}
Resumo final e possíveis trabalhos futuros.

% Exemplo de figura
% \section{Exemplo de Figura}
% \begin{figure}[h!]
% 	\centering
% 	\includegraphics[width=0.5\textwidth]{exemplo.png}
% 	\caption{Exemplo de imagem}
% 	\label{fig:exemplo}
% \end{figure}

% Referências
% \begin{thebibliography}{9}
% 	\bibitem{ref1} Autor. \textit{Título do livro/artigo}. Editora, ano.
% \end{thebibliography}

\end{document}
